\documentclass[11pt,a4paper]{article}

\usepackage{fontspec}
\setmainfont{Arial}[Path=/System/Library/Fonts/Supplemental/, Extension=.ttf, UprightFont=Arial, BoldFont=Arial Bold, ItalicFont=Arial Italic, BoldItalicFont=Arial Bold Italic]
\usepackage[bulgarian]{babel}
\usepackage[a4paper,margin=2.2cm]{geometry}
\usepackage{amsmath}
\usepackage{booktabs}
\usepackage{array}
\usepackage{tabularx}
\usepackage{enumitem}
\usepackage{xcolor}
\usepackage{hyperref}
\usepackage{fancyhdr}

\definecolor{navy}{HTML}{1F3A5F}
\definecolor{todo}{HTML}{B23A48}

\hypersetup{colorlinks=true, linkcolor=navy, urlcolor=navy}

\newcommand{\todo}[1]{\textcolor{todo}{[#1]}}

\setlength{\parindent}{0pt}
\setlength{\parskip}{4pt}
\setlist{nosep}

\pagestyle{fancy}
\fancyhf{}
\fancyhead[L]{\small Курсов проект: STEM среда в СМГ}
\fancyhead[R]{\small Управление на проекти}
\fancyfoot[C]{\thepage}

\renewcommand{\headrulewidth}{0.4pt}

\begin{document}

\begin{center}
  {\LARGE\bfseries Курсова работа / Курсов проект\par}
  \vspace{4pt}
  {\Large за разработване и управление на интердисциплинарен STEM център\\с интеграция на изкуствен интелект\par}
  \vspace{14pt}
  {\large Училище: Софийска математическа гимназия ``Паисий Хилендарски'' (СМГ)\par}
  {\large Учебна дисциплина: Управление на проекти \quad 2026/2027 г.\par}
  \vspace{8pt}
  {\normalsize Екип: \todo{имена на двамата студенти}\par}
\end{center}

\vspace{6pt}
\noindent\textbf{Основен въпрос:} \emph{``Как чрез оптимална комбинация от технологии, учебни дисциплини, данни, физическа среда и изкуствен интелект може да бъде създаден ефективен интердисциплинарен STEM център?''}

\medskip
\noindent\textbf{Отговор в една кратка теза.} Предложеното решение е по-интегрирано, по-ефективно и по-устойчиво от отделни лаборатории по предмети, защото: (1) едни и същи сензори, контролери и данни се използват в поне четири предмета; (2) една и съща инфраструктура обслужва три направления и класни стаи (ВОСКС); (3) AI се използва за анализ и управление на средата, а не като изолирана функция; (4) разходите се контролират чрез проверка на всеки ред срещу публична цена с източник.

\section{Дефиниране на проекта и SMART цели}

\subsection{Описание на идеята за изграждане на STEM среда}

Училищният STEM център на СМГ се изгражда като една интегрирана среда с три технологични направления:
\begin{itemize}
  \item \textbf{Направление 1 – Математика и информатика:} програмиране на микроконтролери и едноплатни компютри, анализ на данни, прости модели за машинно обучение, графично калкулиране.
  \item \textbf{Направление 2 – Роботика и кибер-физични системи:} IoT устройства, колесни роботи, роботизирана ръка, учебни дронове, прототипиране на 3D печат.
  \item \textbf{Направление 3 – Природни науки:} лабораторни сензори (pH, EC, температура, влажност), интелигентна оранжерия (Smart Greenhouse), метеорологична станция, микроскопия, лабораторни измервания.
\end{itemize}

Към трите направления се добавя \textbf{ВОСКС} – 8 високотехнологично оборудвани класни стаи с интерактивен дисплей, документна камера, безжично представяне и учителски компютър.

\textbf{Специализирано оборудване.} Предвижда се ограничено допълнително специализирано оборудване (до 20\% от разходите за оборудване), само за Направление 3 – напр. калибровъчни буфери за pH и лабораторна везна с калибровъчна маса. Останалото оборудване е стандартно и купувано от публични магазини.

\textbf{Пространство.} Гъвкави мебели и подвижни работни маси за групова работа, отделна зона за презентации и място за учителите. Конкретният план за СМР се изготвя след оглед на помещенията \todo{изисква оглед и оферта от строителна фирма}.

\subsection{Учебни предмети, свързани с дейността на центъра}

Минимум четири учебни предмета са включени в интегрираното преподаване:
\begin{enumerate}
  \item \textbf{Информатика и информационни технологии} – програмиране на контролери, обработка на данни, обучение на AI модели. Връзка: Направление 1 и 2.
  \item \textbf{Математика} – статистика, моделиране, оптимизация на параметри, анализ на регресия и прогнози. Връзка: Направление 1 и 3.
  \item \textbf{Физика} – електрически вериги, сензори, механика на роботи и дронове, оптика на микроскоп. Връзка: Направление 2 и 3.
  \item \textbf{Биология и здравно образование} – растеж на растенията, условия на средата, наблюдения и измервания. Връзка: Направление 3.
\end{enumerate}

Защо са избрани: тези четири предмета покриват измерванията (физика, биология), обработката им (математика) и програмирането (информатика), което позволява един и същ сензорен поток да се използва в поне три предмета едновременно.

Как ще се преподават интегрирано: учебен сценарий "Умна оранжерия" – ученици от биология поставят хипотеза за влажността на почвата, физиката настройва сензорите и помпата, информатиката програмира контролера, а математиката анализира данните и построява прогнозен модел.

\subsection{SMART цели}

\begin{description}[leftmargin=3cm, style=nextline]
  \item[S -- Specific] Да се изгради STEM среда с три направления и ВОСКС в СМГ, която решава проблема с изолираните лаборатории и ниското използване на оборудването.
  \item[M -- Measurable] (1) Минимум 4 интегрирани учебни предмета. (2) Минимум 6 интердисциплинарни STEM сценария за учебната година. (3) Минимум 70\% заетост на оборудването, отчетена по график на ползване. (4) Минимум 2 приложения на AI в обучението и управлението. (5) Минимум 200 ученици годишно работят в средата по график. Методология: логове от оборудването, график на ползване, анкети за учители и ученици в края на всеки семестър.
  \item[A -- Achievable] Финансиране в рамките на лимита за STEM средства (600\,000 лв. с ДДС) и ВОСКС (80\,000 лв. с ДДС, минимум 8 оборудвани класни стаи). Обучение на 6 учители през първата година. Определени отговорници: ръководител на проекта и координатор по STEM.
  \item[R -- Relevant] Развива STEM компетентности и интердисциплинарно мислене; свързва учебното съдържание с реални данни и проекти (IoT, анализ на данни, AI).
  \item[T -- Time-bound] Март – август 2027: проектиране и СМР. Септември 2027: въвеждане в експлоатация и първи сценарии. Юни 2028: оценка на резултатите с модела STEM Integration Score (раздел 8).
\end{description}

\section{Модел за интердисциплинарна STEM интеграция}

Моделът свързва учебни дисциплини, STEM области, технологии, оборудване, данни, AI и учебни дейности (\emph{STEM Integration Matrix}).

\begin{table}[h]
\centering
\small
\begin{tabularx}{\textwidth}{@{}>{\bfseries}l X X X X@{}}
\toprule
Учебна дисциплина & STEM област & Технология & Оборудване & Данни \,/\, AI \,/\, учебна дейност \\
\midrule
Биология & природни науки & IoT сензори & Smart Greenhouse, сензор за влажност & температура, влажност \,/\, AI прогноза \,/\, прогноза за условията за растеж \\
Физика & физика и инженерство & ел. вериги, микроконтролер & Микроконтролер, сензорен комплект & напрежение, ток \,/\, автоматично управление \,/\, изграждане на помпена система \\
Информатика & програмиране & Python, Arduino C & Едноплатен компютър, микроконтролер & сензорни данни \,/\, ML модел \,/\, автоматично управление \\
Математика & анализ на данни & облачна платформа, статистика & Преносим компютър & исторически данни \,/\, AI регресия \,/\, статистически и предиктивен анализ \\
\bottomrule
\end{tabularx}
\end{table}

\textbf{Пример за верига (Биология $\rightarrow$ Природни науки $\rightarrow$ IoT сензори $\rightarrow$ Smart Greenhouse $\rightarrow$ влажност $\rightarrow$ AI анализ $\rightarrow$ прогноза):} учениците от биология определят оптималните условия за растеж; сензорите мерят влажност и температура; данните се събират в облачна платформа; AI моделът прогнозира кога ще е нужно поливане; решението се потвърждава от учител (human-in-the-loop).

\textbf{Графична схема на връзките.} Един сензорен поток (температура, влажност, светлина) се използва от поне три предмета: физика (калибриране и вериги), биология (условия за растеж) и математика (статистика и прогноза). Същата платформа обслужва ВОСКС класните стаи за визуализация на резултатите.

\section{AI интеграция в STEM центъра}

Предлагаме две реални приложения на AI, като и двете са в различни направления.

\subsection{AI приложение 1: прогноза и управление на Smart Greenhouse (Направление 3)}

\begin{tabularx}{\textwidth}{@{}>{\bfseries}lX@{}}
\toprule
Вид & AI за анализ на данни и AI за управление на STEM средата \\
\midrule
Входни данни (Input) & Почвена влажност, температура, влажност на въздуха, осветеност – от IoT сензори в оранжерията \\
Обработка на данните (Data Processing) & Събиране през Wi-Fi към облачна платформа, почистване на липсващи стойности, агрегиране на 10-минутни интервали, съхранение \\
AI модел или услуга (AI Model/Service) & Модел за прогноза на времевия ред (регресия или LSTM); алтернатива: готов облачен сервис за прогнози \\
Резултат (Output) & Прогноза на почвената влажност за следващите 24 часа и предупреждение при аномалия \\
Решение или действие (Decision/Action) & Препоръка за поливане, която се потвърждава от ученик или учител; автоматично поливане само в рамките на зададен лимит \\
\bottomrule
\end{tabularx}

\subsection{AI приложение 2: AI асистент за индивидуални задачи и обратна връзка по програмиране (Направление 1)}

\begin{tabularx}{\textwidth}{@{}>{\bfseries}lX@{}}
\toprule
Вид & AI за обучение \\
\midrule
Входни данни (Input) & Код на ученика, зададени задачи, резултати от предишни задачи \\
Обработка на данните (Data Processing) & Анонимизиране на имената, филтриране на лични данни, съхранение на резултатите в локална база \\
AI модел или услуга (AI Model/Service) & Езиков модел чрез API (облачна услуга) с инструкции, настроени за учебни цели \\
Резултат (Output) & Генерирани индивидуални задачи с различна трудност и коментари към кода \\
Решение или действие (Decision/Action) & Учителят одобрява задачата и коментара преди да стигнат до ученика \\
\bottomrule
\end{tabularx}

\subsection{Рискове на AI приложенията и мерки за ограничаването им}

\begin{tabularx}{\textwidth}{@{}>{\raggedright\arraybackslash}p{4.2cm}X@{}}
\toprule
Риск & Мярка \\
\midrule
Неточност на прогнозите & Сравнение с реални измервания на всяка седмица; праг на доверие; ръчно потвърждение на поливането \\
Генериране на невярна информация (hallucinations) & Проверка от учител преди публикуване; ограничаване на отговорите до утвърдена база задачи \\
Пристрастие (bias) & Обучение на модела върху данни от поне два сезона и различни условия \\
Защита на данните & Без имена и лични данни в облака; съхранение на локален сървър \\
Киберсигурност & Защитена Wi-Fi мрежа с отделен сегмент за IoT; смяна на паролите по график; актуализации на контролерите \\
Прекомерна зависимост от AI & Задължителен human-in-the-loop; учителите запазват ръчен режим на всяко устройство \\
\bottomrule
\end{tabularx}

\subsection{Тестване на AI интеграцията: цена и време}

Професорът повдигна въпрос за външното тестване на интеграцията. Предлагаме следния план; стойностите са предположения за планиране и ще се заменят с реални данни след тестовете.

\begin{tabularx}{\textwidth}{@{}>{\raggedright\arraybackslash}p{5.2cm}XX@{}}
\toprule
Дейност & Време (оценка) & Разход (оценка) \\
\midrule
Тестов набор от 500 заявки (задачи и код) за AI асистента & 1 седмица, 1 специалист & Предположение: до 40 лв. за API заявки \\
Тест на прогнозния модел върху исторически данни (1 сезон) & 2 седмици, 1 специалист & Предположение: 0 лв. (локален софтуер) \\
Външна проверка на интеграцията от доставчик на облачната платформа & 3--5 работни дни & \todo{оферта от доставчика} \\
Проверка за сигурност (външна) на IoT сегмента & 1 седмица & \todo{оферта от външна фирма} \\
\bottomrule
\end{tabularx}

\section{Бизнес и технологичен анализ -- SWOT}

\begin{tabularx}{\textwidth}{@{}>{\bfseries}lX@{}}
\toprule
Компонент & Анализ за проекта \\
\midrule
Strengths (Силни страни) & Интердисциплинарна интеграция в четири предмета; съвременно оборудване (IoT, AI, роботика); една платформа за данни за всички направления; споделено използване на ресурси. \\
Weaknesses (Слаби страни) & Високи разходи за реализация; нужда от обучение на учителите; техническа сложност при интеграцията на системите; зависимост от поддръжка на оборудването. \\
Opportunities (Възможности) & Развитие на AI, роботика и виртуални лаборатории; нови интердисциплинарни програми; участие в конкурси и партньорства с университети. \\
Threats (Заплахи) & Технологично остаряване на оборудването; киберрискове; недостатъчно използване на оборудването; зависимост от доставчици и промени в цените. \\
\bottomrule
\end{tabularx}

\textbf{Стратегически изводи.} SO: използване на силата на интеграцията за участие в конкурси и партньорства. WO: обучение и споделени сценарии намаляват тежестта на сложността. ST: избор на открито (open-source) софтуер и стандартни контролери намалява риска от остаряване и зависимост от доставчик. WT: поетапно внедряване с оценка на всяка фаза намалява риска от надхвърляне на бюджета.

\section{Управление на риска}

\subsection{Risk Breakdown Structure (RBS)}

\begin{itemize}
  \item \textbf{Технологични рискове:} несъвместимост на оборудването, технически неизправности, ограничения на технологиите.
  \item \textbf{Финансови рискове:} превишаване на бюджета, непредвидени разходи, недостатъчно финансиране.
  \item \textbf{Организационни рискове:} забавяне на дейностите, неефективно разпределение на ресурсите.
  \item \textbf{Образователни рискове:} недостатъчна подготовка на учителите.
  \item \textbf{AI рискове:} неточност, пристрастност, прекомерна зависимост от AI.
  \item \textbf{Рискове, свързани с данните:} загуба или нарушена цялост на данните.
  \item \textbf{Киберсигурност:} неоторизиран достъп, кибератаки.
  \item \textbf{Инфраструктурни рискове:} недостатъчна електрическа или мрежова инфраструктура.
  \item \textbf{Интеграционни рискове:} затруднено свързване на компонентите в единна среда.
\end{itemize}

\subsection{Risk Matrix}

Оценката е качествена и количествена: $\text{Risk Score} = \text{Probability} \times \text{Impact}$ (1--5 всяка). Нива: до 5 -- ниско; 6--12 -- средно; 13 и повече -- високо.

\begin{table}[h]
\centering
\small
\begin{tabularx}{\textwidth}{@{}>{\raggedright\arraybackslash}p{5.6cm}cccX@{}}
\toprule
Риск & P & I & Score & Мярка и отговорно лице \\
\midrule
Превишаване на бюджета (цени, доставки) & 3 & 4 & 12 (средно) & Проверка на всеки ред с източник; резерв 5\%; финансов отговорник \\
Забавяне на СМР и електроинсталация & 3 & 4 & 12 (средно) & Оглед преди договора; поетапно плащане; координатор по СМР \\
Неточност и пристрастност на AI модела & 4 & 3 & 12 (средно) & Седмичен мониторинг; ръчно потвърждение; AI специалист \\
Липса на обучен персонал & 3 & 3 & 9 (средно) & Обучение на 6 учители; учители лидери; координатор по STEM \\
Кибератака към облачната платформа & 2 & 5 & 10 (средно) & Отделна мрежа за IoT; смяна на пароли; външен одит \\
Защита на личните данни & 2 & 4 & 8 (средно) & Анонимизиране; локално съхранение; съгласие на родителите \\
\bottomrule
\end{tabularx}
\end{table}

\textbf{Приоритет:} рискът от превишаване на бюджета и рискът от забавяне на СМР имат най-висок приоритет, защото влияят на всички останали дейности.

\subsection{FMEA (Failure Mode and Effects Analysis)}

За всеки отказ: Impact (I), Probability (P) и Detection (D) по скала 1--10. Risk Value $=$ I $\times$ P $\times$ D.

\begin{table}[h]
\centering
\small
\begin{tabularx}{\textwidth}{@{}>{\raggedright\arraybackslash}p{3.4cm}>{\raggedright\arraybackslash}p{3.6cm}cccc>{\raggedright\arraybackslash}X@{}}
\toprule
Компонент & Отказ & I & P & D & RV & Превантивна мярка \\
\midrule
IoT сензорна система & Отказ на сензор за влажност & 6 & 5 & 3 & 90 & Дублиране на ключовите сензори; калибриране на всеки месец \\
AI модел & Грешна прогноза (drift) & 7 & 4 & 6 & \textbf{168} & Мониторинг на точността; ръчно потвърждение; повторно обучение \\
Smart Greenhouse & Непрекъсната работа на помпата & 8 & 3 & 2 & 48 & Хардуерен таймер; ограничение на обема; аварийно изключване \\
Облачна платформа & Прекъсване на връзката & 5 & 4 & 2 & 40 & Локално кеширане на данните; автоматично повторно изпращане \\
Мрежова инфраструктура & Слабо Wi-Fi покритие в класа & 6 & 4 & 3 & 72 & Оглед на покритието; допълнителна точка за достъп \\
\bottomrule
\end{tabularx}
\end{table}

\textbf{Критичен отказ:} AI моделът (RV = 168), поради което за него е предвидено ръчно потвърждение на всяко действие и седмичен мониторинг.

\section{Управление на проектния екип}

\textit{(Раздел за студенти с курсов проект.)}

\subsection{Организационна структура и роли}

\begin{tabularx}{\textwidth}{@{}>{\raggedright\arraybackslash}p{5cm}X@{}}
\toprule
Роля & Отговорности \\
\midrule
Ръководител на проекта (Project Manager) & Планиране, срокове, ресурси, контрол; води задачите в Jira/таблица \\
Специалист по STEM интеграция & Координира връзките между предмети, направления и оборудване \\
Специалист по AI и машинно обучение (AI/ML Specialist) & AI приложения 1 и 2; тестване на интеграцията \\
Софтуерен разработчик & Интеграция на сензорите, облачната платформа и ВОСКС \\
Специалист по инфраструктура & Оглед, мрежа, електрически изисквания, СМР \\
Анализатор на данни & Събиране и анализ на данните за ефективността на средата \\
\bottomrule
\end{tabularx}

\todo{разпределението на ролите между двамата студенти да се попълни}

\subsection{Развитие на екипа по модела на Тъкман}

\begin{description}[leftmargin=3cm, style=nextline]
  \item[Forming] Определяне на целите, ролите и отговорностите. Управленско действие: писмен списък с роли и срокове.
  \item[Storming] Разногласия по технологичните решения и разпределението на оборудването. Управленско действие: решения се записват, а спорните въпроси се решават с данни (цени, тестове).
  \item[Norming] Правила за работа: общ списък на задачите, срещи веднъж седмично, решения в писмен вид.
  \item[Performing] Изпълнение на дейностите, решаване на проблемите и постигане на целите. Управленско действие: преглед на риска и на графика всяка седмица.
  \item[Adjourning] Финализиране на проекта, оценка на резултатите и анализ на придобития опит. Управленско действие: ретроспекция и доклад.
\end{description}

\todo{да се отбележи текущият етап на екипа и датите на преходите}

\section{Използване на изкуствен интелект при управлението на проекта}

AI е използван като асистент за планиране, структуриране, търсене на данни и подготовка на материали. Всички резултати са проверени от екипа преди употреба.

\subsection{AI Project Log}

\begin{tabularx}{\textwidth}{@{}>{\raggedright\arraybackslash}p{3.4cm}>{\raggedright\arraybackslash}p{3cm}>{\raggedright\arraybackslash}p{3.8cm}X@{}}
\toprule
Задача & Използван AI инструмент & Prompt (заявка) & Генериран резултат \,/\, Приет/отхвърлен \,/\, Проверка и оценка \\
\midrule
Структура и текст на слайдовете за раздел 1 & Claude Code & Създай 2 слайда: идея за STEM среда и SMART цели & Генерирани 2 слайда. Приети с редакции. Проверка: текстът е прочетен; визуалният вид не е проверен на екран. \\
Търсене на цени на оборудване & Claude Code с браузър и помощни агенти & Намери най-евтин магазин за всеки продукт на pazaruvaj.com & Част от продуктите са намерени, други не. Приети само цени, които са прочетени от страницата; отхвърлени са продукти без съвпадение със спецификацията. \\
Структура на LaTeX документа & Claude Code & Изгради задание по образеца от PDF & Приет след компилация и преглед на всеки раздел. \\
\bottomrule
\end{tabularx}

\subsection{Сравнителен анализ: традиционен и AI-подпомогнат подход}

\begin{tabularx}{\textwidth}{@{}>{\raggedright\arraybackslash}p{4.4cm}XX@{}}
\toprule
Показател & Human-only & AI-assisted \\
\midrule
Време за изпълнение (ценово проучване, 20 продукта) & \todo{измерва се} & \todo{измерва се} \\
Качество на резултатите (пълнота, приложимост) & \todo{оценка след измерване} & \todo{оценка след измерване} \\
Точност (грешки, нужни корекции) & \todo{брой грешки} & \todo{брой грешки} \\
Брой генерирани алтернативи & 1--2 & 3--4 \\
\bottomrule
\end{tabularx}

Изводи (предварителни): AI намалява времето за първоначалното търсене и структуриране, но всяка цена и всяко твърдение изисква проверка от човек. Не е измерен резултат, затова изводите са предварителни.

\section{Оптимизация на STEM интеграцията}

\textit{(Раздел за студенти с курсов проект.)}

\subsection{Критерии, тегла и оценки}

Всеки критерий се оценява по скала 1--5. Теглата се определят според значимостта за проекта и сумата им е 1,00.

\begin{tabularx}{\textwidth}{@{}>{\raggedright\arraybackslash}p{5.2cm}cXX@{}}
\toprule
Критерий & Тегло ($W_i$) & Пояснение & Предложение (A) \\
\midrule
Интеграция на учебни дисциплини & 0{,}25 & брой предмети, свързани с един сценарий & 5 \\
Споделено използване на технологии & 0{,}20 & едни и същи сензори и платформа за няколко предмета & 5 \\
Степен на използване на оборудването & 0{,}15 & очаквана заетост по график & 4 \\
AI интеграция & 0{,}15 & брой AI приложения и връзки с учебни дейности & 5 \\
Взаимосвързани информационни потоци & 0{,}10 & връзка между сензори, облак, ВОСКС & 5 \\
Разходи за реализация & 0{,}10 & по-ниска цена = по-висока оценка & 3 \\
Сложност на управлението & 0{,}05 & по-проста поддръжка = по-висока оценка & 2 \\
\midrule
Общо & 1{,}00 & & \\
\bottomrule
\end{tabularx}

\subsection{Сравнение на три варианта}

\begin{tabularx}{\textwidth}{@{}>{\raggedright\arraybackslash}p{4.6cm}cccc@{}}
\toprule
Критерий (тегло) & A: интегрирана среда (предложение) & B: отделни лаборатории по предмети & C: без ВОСКС, две направления \\
\midrule
Интеграция на учебни дисциплини (0{,}25) & 5 & 2 & 4 \\
Споделено използване на технологии (0{,}20) & 5 & 2 & 3 \\
Степен на използване на оборудването (0{,}15) & 4 & 2 & 4 \\
AI интеграция (0{,}15) & 5 & 1 & 3 \\
Взаимосвързани информационни потоци (0{,}10) & 5 & 1 & 3 \\
Разходи за реализация (0{,}10) & 3 & 4 & 4 \\
Сложност на управлението (0{,}05) & 2 & 4 & 3 \\
\midrule
\textbf{STEM Integration Score} $=\sum W_i X_i$ & \textbf{4,50} & \textbf{2,05} & \textbf{3,50} \\
В проценти от максимума (5,00) & 90\% & 41\% & 70\% \\
\bottomrule
\end{tabularx}

\textbf{Извод.} Вариант A има най-висок резултат (4,50 от 5,00). Той е избран, защото печели при интеграцията, AI и информационните потоци, които са основата на проекта; слабостта му са разходите и сложността, които се компенсират с поетапно внедряване. Оценките са на екипа и се проверяват при всяка фаза.

\section{Финансов анализ}

Изходният резултат е Excel файлът \texttt{concept-stem-smart-finance.xlsx} (в същата папка). Той съдържа: ред по ред оборудване, софтуер, обзавеждане и СМР с количество, единична цена без ДДС, общо с ДДС и източник на цената (връзка към магазин в pazaruvaj.com); проверка на лимитите за STEM средства и ВОСКС.

\begin{tabularx}{\textwidth}{@{}Xr@{}}
\toprule
Раздел на бюджета (в лв. с ДДС) & Сума \\
\midrule
Лимит за STEM средства (СМГ) & 600\,000 \\
Оборудване (направления) -- минимум 45\% (сумата е 100\% от заявката) & 59\,887{,}74 \\
Обзавеждане и СМР -- максимум 45\% & 0 \\
Специализирано оборудване -- максимум 20\% от оборудването & 0 \\
Софтуер -- максимум 10\% от STEM средствата & 0 \\
Общо заявка за STEM среда & 59\,887{,}74 \\
Лимит за ВОСКС (минимум 8 класни стаи) & 80\,000 \\
Заявка за ВОСКС (8 класни стаи, оборудване) & 44\,371{,}45 \\
Общо STEM и ВОСКС & 104\,259{,}19 \\
Организация и управление (1\%) & 1\,042{,}59 \\
Общо за изграждане на STEM среда & 105\,301{,}78 \\
\bottomrule
\end{tabularx}

\textbf{Методология.} Всяка позиция има две оферти от различни магазини (от pazaruvaj.com или директно от магазина), с цена и линк до всяка; в лист \texttt{concept-stem-smart-finance.xlsx} единичната цена е средната стойност на двете. Цените са към 05.11.2026. Цените в евро са превърнати в лева по фиксирания курс 1{,}95583; ДДС е 20\%. Доставката не е включена, освен ако не е посочено. Линковете към pazaruvaj са към страницата на продукта, а не към директната страница на магазина, защото пренасочващите връзки са с времеви ключ. Някои позиции нямат втори магазин и са отбелязани в Excel.

\textbf{Защо заявката е под лимита.} Прегледани са 29 позиции: 18 са намерени с цена (част от тях са частични съвпадения, отбелязани в Excel), а 11 не са намерени и са с количество 0 в Excel. Обзавеждането и СМР не са включени, защото за тях няма оферта; без тях заявката не може да се сравни с 600\,000 лв. Преди подаване ще се добавят офертите за СМР и обзавеждане, като оборудването остава поне 45\% от общата сума. Липсващите AI ускорител, сензорни модули за pH и EC, смарт оранжерия, роботизирана ръка и учебен дрон също трябва да се намерят от други търговци.

\section{Финална презентация}

Продължителност: 10--12 минути. Основният акцент е върху ефективната интеграция между учебните дисциплини, технологиите, информационните системи, данните, изкуствения интелект и проектните дейности.

\end{document}
